%===============================================================================
% Zweck:    KTR-Seminar-Vorlage
% Erstellt: 16.10.2007
% Updated:  15.04.2013
% Autor:    U.K. / M.G.
%===============================================================================

%===============================================================================
% Zum Kompilieren pdflatex und bibtex ausführen.
%	Konfiguration: in texmaker unter Benutzer -> Benutzerbefehle einen neuen Befehl anlegen: pdflatex -synctex=1 -interaction=nonstopmode %.tex | bibtex % | makeindex %.nlo -s nomencl.ist -o %.nls | pdflatex -synctex=1 -interaction=nonstopmode %.tex | pdflatex -synctex=1 -interaction=nonstopmode %.tex
%	Entsprechende Informationen in den config/metainfo verändern
% Zur Auswahl der Sprache im folgenden Befehl
% ngerman für deutsch eintragen, english für Englisch.
%===============================================================================

% Options ngerman, english
\newcommand{\lang}{english}

\documentclass[pdftex, journal, onecolumn, a4paper, 12pt, \lang]{IEEEtran}

%===============================================================================
% zentrale Layout-Angaben und Befehle
%===============================================================================
% What You should change:
% Here goes your name
\author{Frank Kaiser}
% and the title of your seminar
\newcommand{\subtitle}{Fog Computing}
% the date of the submission
\date{\today}

% What is already done
\newcommand{\semester}{\ifthenelse{\equal{\lang}{english}}{Sommersemester 2017}{Summer Term 2017}}
\title{\ifthenelse{\equal{\lang}{ngerman}}{Moderne Internet-Dienste in Software-definierten und\\ informationszentrierten Netzen der\\[.5em] n\"achsten Generation}{%
Modern Internet Services in Software Defined and\\[.5em] Information Centric Networks of the next Generation
}}
\newcommand{\supervisor}{Prof. Dr. Udo Krieger}
%\institute[shortinstitute]{Institute}

%\titlegraphic{\includegraphics[width=13mm,height=13mm]{image/logo}}
%\subject{subject}
%\keywords{keywords}
%===============================================================================
% Zweck: KTR-Seminar-Vorlage in Anlehung an G. Wirtz, Lehrstuhl Praktische Informatik
%===============================================================================
%===============================================================================
% zentrale Layout-Angaben und Befehle
%===============================================================================
%
% \usepackage{german}
\usepackage{ifthen}
% Language Selection for Babel
\ifthenelse{\equal{\lang}{ngerman}}{\usepackage[german, \lang]{babel}}{\usepackage[\lang]{babel}}

\usepackage[utf8]{inputenc}
\usepackage{fancyhdr}
\usepackage[T1]{fontenc}
\usepackage{ae}
\usepackage{color}
\usepackage{amsmath}
\usepackage{amsfonts}
%

\usepackage{caption}

%%   Fuer anspruchsvolle Tabellen   %%
\usepackage{longtable, colortbl}
\usepackage{multicol, multirow}

\usepackage[pdftitle={\@title},pdfauthor={\@author},pdftex,bookmarksopen,bookmarksnumbered,pdfborder=0]{hyperref}
\usepackage[pdftex]{graphicx}
\usepackage{float}
\usepackage{tikz}
\usepackage{pgfplots}
\usetikzlibrary{arrows,shapes,fit,positioning,snakes,backgrounds,shadows}

\pdfcompresslevel=9

% Code-Hervorhebung
% Quellcode
\usepackage{verbatim}            % Quellcode einbinden (\verbatiminput) standardpaket
\usepackage{moreverb} 
% PseudoCode
\usepackage{algorithm}
\usepackage{algpseudocode}
%\usepackage{algorithmicx}
\floatname{algorithm}{Algorithmus}
\algrenewcommand{\algorithmiccomment}[1]{\hskip1em\textcolor{gray!60}{$\rhd$ #1}}
\ifthenelse{\equal{\lang}{ngerman}}{%
\renewcommand{\listalgorithmname}{Algorithmen}
\def\algorithmautorefname{Algorithmus}
}{%
\renewcommand{\listalgorithmname}{List of Algorithms}
\def\algorithmautorefname{Algorithm}
}


%%   intoc zur Aufnhame des Abkuerzungs- und Symbolverzeichnisses ins Inhaltsverzeichnis  
\usepackage[intoc]{nomencl}
\setlength{\nomlabelwidth}{.20\hsize}
%\renewcommand{\nomlabel}[1]{#1 \dotfill}
\setlength{\nomitemsep}{-\parsep}
\makenomenclature
\ifthenelse{\equal{\lang}{ngerman}}{%
\renewcommand{\nomname}{Abk\"urzungsverzeichnis}
%\newcommand{\nomaltname}{Symbolverzeichnis}
}{%
\renewcommand{\nomname}{List of Abbreviations}
%\newcommand{\nomaltname}{List of Symbols}
}%
%\newcommand{\nomaltpreamble}{}
%\newcommand{\nomaltpostamble}{}
%\newcommand{\usetwonomenclatures}{\nomenclature[\switchnomitem]{}{}}
%\newcommand{\switchnomitem}{R}
%\renewcommand{\nomgroup}[1]{%
%\ifthenelse{\equal{#1}{\switchnomitem}}{\switchnomalt}{}}
%\newcommand{\switchnomalt}{%
%\end{thenomenclature}
%\newpage
%\renewcommand{\nomname}{\nomaltname}
%\renewcommand{\nompreamble}{\nomaltpreamble}
%\renewcommand{\nompostamble}{\nomaltpostamble}
%\begin{thenomenclature}
%}


%%   Hervorhebung der Abkuerzungsbuchstaben   %%
\usepackage[normalem]{ulem}
\newcommand{\m}[1]{\uline{#1}}

%\usepackage[dvips]{rotating}
%
% ausf\"{u}hrlichere Fehlermeldungen
\errorcontextlines=999
%
% Page-Layout: A4 aus Header
% Alternative
\setlength\headheight{14pt}
\setlength\topmargin{-15,4mm}
\setlength\oddsidemargin{-0,4mm}
\setlength\evensidemargin{-0,4mm}
\setlength\textwidth{160mm}
\setlength\textheight{252mm}
%
%% Absatzeinstellungen
\setlength\parindent{0mm}
\setlength\parskip{2ex}


\makeatletter
\renewcommand{\maketitle} {
  \begin{titlepage}
  \centering
    \begin{minipage}[t]{16cm}
      \hfill
      \begin{minipage}{12cm}
      \ifthenelse{\equal{\lang}{ngerman}}{%
        \centering
        Otto-Friedrich-Universit\"at Bamberg
        \\[12pt]%
        {\Large Professur f\"ur Informatik\\[.5em]%
        \large insbesondere Kommunikationsdienste,\\[.5em]%
        Telekommunikationsdienste und Rechnernetze}%
        }{%
        \centering
        Otto-Friedrich-University of Bamberg
        \\[12pt]
        {\Large Professorship for Applied Computer Science,\\[.5em]%
        \large Cognitive Systems}%
        }
      \end{minipage}
      \hfill
      \begin{minipage}{3cm}
        \includegraphics[height=28mm]{config/images/logo} %height=26mm
      \end{minipage}
    \end{minipage}\\[70pt]%[50pt]
   % 
    {\Large\bf \ifthenelse{\equal{\lang}{ngerman}}{Ausarbeitung des KTR-Seminars}{Seminar on}}
    \\[36pt]
   % oder: {\LARGE Ausarbeitung des KTR-Seminars}\\[12pt]
    {\LARGE \@title}\\[80pt]
    {\Large\bf \ifthenelse{\equal{\lang}{ngerman}}{Thema:}{Topic:}}\\[36pt]
    {\LARGE\bf \subtitle}\\
    \vfill
    \begin{minipage}{\textwidth}
      \center
      \ifthenelse{\equal{\lang}{ngerman}}{Vorgelegt von:}{Submitted by:}\\
      {\Large \@author \\[18pt]}
      \ifthenelse{\equal{\lang}{ngerman}}{Betreuer:}{Supervisor:} \supervisor \\[12pt]
      Bamberg, \@date\\
      \semester
    \end{minipage}
  \end{titlepage}
}
\makeatother
%
\usepackage[intoc]{nomencl}
\ifthenelse{\equal{\lang}{ngerman}}{%
\renewcommand{\nomname}{Abkürzungsverzeichnis}}%
{\renewcommand{\nomname}{List of Abbreviations}}
\renewcommand{\nomlabel}[1]{#1 \dotfill}
\makenomenclature

% Einbindung eines Bildes mit angegebener Breite
% #1 = label f\"{u}r \ref-Verweise
% #2 = Name des Bildes ohne Endung relativ zu Bilder-Verzeichnis
% #3 = Beschriftung
% #4 = Breite des Bildes im Dokument in cm
\newcommand{\bildw}[4]{%
  \begin{figure}[htb]%
    \centering
    \includegraphics[width=#4cm]{Bilder/#2}%
    \vskip -0.3cm%
    \caption{#3}%
    \vskip -0,2cm%
    \label{#1}%
  \end{figure}%
}
%
% Einbindung eines Bildes mit Seitenbreite
% #1 = label f\"{u}r \ref-Verweise
% #2 = Name des Bildes ohne Endung relativ zu Bilder-Verzeichnis
% #3 = Beschriftung
\newcommand{\bild}[3]{%
  \begin{figure}[htb]%
    \centering%
    \includegraphics[width=\textwidth]{Bilder/#2}%
    \vskip -0.3cm%
    \caption{#3}%
    \vskip -0,2cm%
    \label{#1}%
  \end{figure}%
}
%
\numberwithin{equation}{section}
%
%===============================================================================
% zentrale Layout-Angaben und Befehle
%===============================================================================
%
%#1 Breite
%#2 Datei (liegt im image Verzeichnis)
%#3 Beschriftung
%#4 Label fuer Referenzierung
\newcommand{\image}[4]{%
\begin{figure}[H]%
\centering%
\includegraphics[width=#1]{image/#2}%
\caption{#3}%
\label{#4}%
\end{figure}%
}

%#1 Datei (liegt im graphic Verzeichnis)
%#2 Beschriftung
%#3 Label fuer Referenzierung
%#4 Skalierungsfaktor
\newcommand{\scaletikzimage}[4]{%
\begin{figure}[H]%
\centering%
\scalebox{#4}{%
\input{graphic/#1.tikz}}%
\caption{#2}%
\label{#3}%
\end{figure}
}

%#1 algorithm name
%#2 algorithm label
%#3 file name in code-folder
\newcommand{\pseudo}[3]{%
\small%
\begin{algorithm}[H]%
\caption{#1}%
\label{#2}%
\input{code/#3.tex}%
\end{algorithm}%
\normalsize%
}


%===============================================================================
% LATEX-Dokument
%===============================================================================

\hyphenation{op-tical net-works semi-conduc-tor}
\begin{document}

\maketitle

\pagenumbering{Roman}
\setcounter{page}{2}
%
\tableofcontents
% Einstellungen f\"{u}r Literaturverzeichnis
\newpage
\addcontentsline{toc}{section}{\listfigurename}
\listoffigures
\newpage
\addcontentsline{toc}{section}{\listtablename}
\listoftables
\newpage
\printnomenclature
%===============================================================================
% LATEX-Dokument: Kapitel laden
%===============================================================================
%
\newpage
\pagenumbering{arabic}
\setcounter{page}{1}


%
% hier einzelne Kapitel mit \input{Kapitel-File} einf\"{u}gen
%
\begin{abstract}
In this paper, we analyze the different ways to disclose software vulnerabilities. Based on our reviewed literature, we found that the hybrid disclosure method is most often the preferred one as it benefits society and software vendors the most, as it's goal is to combine the advantages of other disclosure policies. Additionally we present a short overview on how to disclose a vulnerability based on the CERT/CC disclosure guide. We then continue to show differences in patch behavior between open source and proprietary software vendors and show the risks and incentives when disclosing vulnerabilities. Finally, we give examples of recent vulnerabilities and their disclosures and analyze them according to the findings of this paper.
\end{abstract}

\section{Introduction}

The word vulnerability origins from the Latin word vulnus, or ``wound''. If a human is wounded, it is quite obvious to see a doctor, but if a computer is wounded, nobody might even notice. Therefore, whenever someone encounters a ``wound'', they should disclose it to the corresponding authorities to provide a patch. Similar to our health care system, where many parties need to work together to ensure the patient's health, it is important that software vendors cooperate with researchers, academics, penetration testers or any user of their product to improve their software's security.

%Everyday more and more software based devices are built and distributed in our environment. While they provide useful services or are just very convenient, every one could contain a flaw, allowing adversaries to break into such devices and steal valuable, private information. Therefore it is of greatest importance that software vendors work together with researchers, academics, penetration testers or any user of their product to improve its software security. Since computers became dominant in our lives, different ways to disclose vulnerabilities have emerged. 

In this paper we introduce three different types of disclosure (\autoref{sec:types_disclosure}) after we establish the terminology in \autoref{sec:vuln_background}. In \autoref{sec:guide} we present a guide for coordinated vulnerability disclosure from the perspective of the discoverer of the vulnerability. We proceed with a short comparison between open source and closed source software vendors in terms of disclosure (\autoref{sec:open_vs_prop}) and provide several incentives and risks for disclosure in \autoref{sec:risks_incentives}. Lastly we give some examples of vulnerability disclosure happened in recent times in \autoref{sec:examples}.


%\vspace{50pt}
%\todo{possible quotes for intro?}

%"Vulnerability disclosure is a process through which vendors and vulnerability finders may work
%cooperatively in finding solutions that reduce the risks associated with a vulnerability. It encompasses
%actions such as reporting, coordinating, and publishing information about a vulnerability and its
%resolution." \cite{ISO29147:2014}

%\vspace{50pt}

%Reports that say that something hasn’t happened are always interesting to me, because as we know, there are known knowns; there are things we know we know. We also know there are known unknowns; that is to say we know there are some things we do not know. But there are also unknown unknowns – the ones we don’t know we don’t know. – Donald Rumsfeld

% What is Vulnerability / Bug? Lifecycle Schryen 2011

\newpage
\section{Background}
\subsection{Reinforcement Learning}




\subsection{Deep Learning}\label{sec:deep_learn}


\subsection{Glossary}


\newpage
\section{Fog and IoT}

\subsection{Rising Challenges of IoT}
The Internet of Things is a very promising but also very criticised field.
Approximately 17 - 30 billion IoT devices are assumed to be online by 2020\footnote{http://spectrum.ieee.org/tech-talk/telecom/internet/popular-internet-of-things-forecast-of-50-billion-devices-by-2020-is-outdated}, in cases where it makes sense and where it doesn't. While for investors this can be a huge opportunity, these devices open a very big attack surface and implementing the fog architecture can make things a lot easier. Chiang is naming 5 different main challenges \cite[p. 855 ff.]{DBLP:journals/iotj/ChiangZ16}:

\begin{itemize}
	\item [A] Stringent Latency Requirements: Control systems require latencies between sensor and control node to be within few milliseconds. Some even within tens of a millisecond.
	\item [B] Network Bandwidth Constraints: Large amount of data generated are generated by sensors(e.g.: up to one gigabyte per second for an autonomous car) causing requiring high bandwidth if sent to the cloud.
	\item [C] Resource-Constrained Devices: IoT devices and micro-controllers often have only the resources they need to work. Everyone of them communicating directly with the cloud will cause a lot of overhead.
	\item [D] Cyber Physical Systems: Cyber-physical systems are often required to be online all the time, especially in safety-critical areas.
	\item [E] Uninterrupted Services With Intermittent Connectivity to the Cloud: Functionality of a system has to be working even without network connectivity.
    \item [F] New Security Challenges:
	\begin{itemize}
		\item[1)] Keeping Security Credentials and Software up to Date on Large Number of Devices.
		\item[2)] Protecting Resource-Constrained Devices.
		\item[3)] Assessing the Security Status of Large Distributed Systems in Trustworthy Manner.
		\item[4)] Responding to Security Compromises Without Causing Intolerable Disruptions.
	\end{itemize}
\end{itemize}

\subsection{How Fog overcomes these challenges [working title]}\label{sec:fog_challenges}

The main advantages of the emergent fog architecture are often illustrated as CEAL \cite[p. 858]{DBLP:journals/iotj/ChiangZ16} and \cite[p. 7]{OpenFog}:

\begin{itemize}
	\item Cognition: A fog node is aware of its environment to react fast to events and if needed redistribute the workload in an optimal manner.
	\item Efficiency: Dynamically using and redistributing unused resources from end user devices
	\item Agility: Rapid innovation and affordable scaling.
	\item Latency: Real-time processing and cyber-physical system control.
\end{itemize}

\subsubsection{Dynamic discovery}
The Internet of Things is a very heterogeneous field. There is a large variety in sensors, that output  different kinds of data over different kinds of communication protocols. That's why a adequate architecture needs to be capable of dynamic discovery and connection of edge devices while providing a high security standard.\\
The discovering of edge nodes can function in two different ways (see )\autoref{img:2_way_discover}):
\begin{itemize}
	\item By letting the edge nodes search continuously for fog gateways which wait for a connection request
	\item By letting the fog gateway search for edge nodes in range and leave the edge devices discoverable
\end{itemize}

\image{8cm}{2_way_discovery.PNG}{Two ways of discovering edge nodes\cite[p. 16]{Perera:2017:FCS:3101309.3057266}}{img:2_way_discover}

Both methods usually use low power communication protocols which will be described in detail in section NOCH EINFÜGEN

% In order to achieve this two things need to happen. First, the edge devices need to be connected to a fog node. Second, the fog node needs to be connected to a IoT cloud platform.\cite[p. 15]{Perera:2017:FCS:3101309.3057266}

\subsubsection{Latency}

Due to evaluating and gathering data closer to the end user, systems can react faster on changing circumstances. This way information does not need to travel to a distant server and back, saving bandwidth and time.

\subsubsection{Network Bandwidth}

Fog controller nodes are capable of distributing the workload along the edge node to cloud continuum in a optimal manner, resulting in less data being sent over the network.


\subsubsection{Uninterrupted Service}

A Fog node is independent of the Internet. It is still able to get data or distribute tasks to its connected edge devices while being disconnected from the rest of the world. This allows local services to stay up even if the node is disconnected.

\subsubsection{Security}

To ensure secure connections from an edge node to the cloud all "things" must "employ a hardware-based immutable root of trust."\cite[p.10]{OpenFog} i.e. to ensure security the hardware itself needs to be protected from unauthorized manipulation. Then on top of that software agents running throughout the infrastructure must be able to ensure the integrity of the root of trust. Due to the nature of edge nodes in being close to the end user, fog nodes often act as first node of access control and encryption enabling the user to decide what privacy-sensitive data should be aggregated before it leaves the node.\cite[p.10]{OpenFog}
\newpage
\section{OpenFog Architecture}

\subsection{The OpenFog Consortium}

The OpenFog Consortium is a composite of multiple big players founded in 2015 in order to solve challenges the usual and wide-spread, centralized cloud architecture could not handle.
The founding members were Arm, Cisco, Dell, Intel, Microsoft and the Princeton University. Right now there are 55 members and the consortium is getting larger as more organizations realize the potential of the Fog Architecture.

\subsection{Fog as a Service}

\image{8cm}{openfog_infrastructure}{OpenFog Infrastacture View\cite[p. 8]{OpenFog}}{img:openfog_infrastracture}

%An OpenFog Fabric is consisting of nodes or layers and may be either centralized or distributed and may be implemented on dedicated hardware, software or both. No matter what of these possibilities is chosen, such fabric distributes resources across available devices, systems and clouds in order to achieve the goal while fulfilling all requirements.\cite[p. 7]{OpenFog}

As stated in Section~\ref{sec:def_fog} Fog's goal is to extend existing cloud architectures and not to replace them. Such Cloud infrastructures are usually offered as Platform as a Service (PaaS).
%In a PaaS you differ between a Service and a Fabric.
%Microsoft, on of the consortium founders, describes it like this for their Azure Cloud Framework documentation\footnote{https://docs.microsoft.com/en-us/azure/service-fabric/service-fabric-cloud-services-migration-differences [03.10.2017]}:

%Cloud Services: Deploying applications as VMs 

%Service Fabric: Deploying applications to existing VMs or machines running Service Fabric on Windows or Linux.

%Analogously OpenFog Fabric and Services are definded:

\textbf{OpenFog Fabric} consists of nodes or layers and may be either centralized or distributed and may be implemented on dedicated hardware, software or both. No matter what of these possibilities is chosen, such fabric distributes resources across available devices, systems and clouds in order to achieve the goal while fulfilling all requirements.\cite[p. 7]{OpenFog}

\textbf{OpenFog Services} are built upon the OpenFog fabric infrastructure. Example services are network accerleration, NFV, SDN, content delivery, device management, device topology, complex event processing, video encoding, field gateway, protocol bridging, traffic offloading, crypto, compression, analytics algorithms/libraries etc.\cite[p. 8]{OpenFog}

\textbf{Devices/Applications} are sensors, microcontrollers, IoT devices running standalone, within a fog deployment or spanning fog deployments.\cite[p. 8]{OpenFog}
 
\textbf{CloudServices} are services that use the cloud for computations and/or work with a large data set or pre-processed edge data.\cite[p. 8]{OpenFog}

\textbf{Security} is important throughout all levels of the architecture as indicated through the vertical pillar in \autoref{img:openfog_infrastracture}. Every unit in every layer needs to participate in state of the art information security practices.\cite[p. 8]{OpenFog}

\textbf{DevOps} (compund word for 'development' and operations') is a software engineering practice which main characteristic is to strongly advocate automation at all steps of software construction. In the OpenFog architecture an efficient set of standard DevOps processes and frameworks enable automation, providing the agility of software upgrades through controlled integration processes.\cite[p. 9]{OpenFog}
\newpage
\section{Conclusion}

Fog is here. And it is here to stay. Fog is just at its start and as the cloud grew over the last 20 years, fog will follow suit and change over the course of its lifespan. Though the basic mantra to let computation, storage and control occur closer to the end user will stay until new drastic technical progress will be made.


\section{Acknowledgement}
%\section{Einleitung}\label{sec:ein}
Einleitung nach \autoref{sec:ein}

\section{Hauptteil}\label{sec:haupt}
\subsection{Bilder und Grafiken}\label{subsec:grafiken}
\subsubsection{Bilder}\label{subsubsec:bilder}
Bilder befinden sich im Image-Orgner und lassen sich mit \textbackslash image\{Breite\}\{Datei im Image-Verzeichnis\}\{Beschriftung\}\{Label\} einbinden. \image{3cm}{logo.png}{Uni-Logo}{img:uni} Die Referenzierung erfolg mittels \textbackslash autoref\{Label\}, also z.B. \autoref{img:uni}.
\subsubsection{Grafiken mit TikZ}
Grafiken im TikZ-Framework\footnote{\url{http://www.tn-home.de/TUGDD/Stuff/TikZ_final.pdf}} lassen sich mit dem Befehl \textbackslash scaletikzimage\{Datei im Image Verzeichnis\}\{Beschriftung\}\{Label\}\{Skalierungsfaktor\} einbinden. \scaletikzimage{tikz}{TikZ-Grafik}{img:tikz}{0.9}
\subsection{Tabellen}
Tabellen lassen sich mit dem Environment\\
\textbackslash begin\{longtable\}[H h t b c]\{Spaltendefinitionen\} ...\\
\qquad\qquad \textbackslash caption\{Tabellenunterschrift\}\\
\qquad\qquad \textbackslash label\{Label\}\\
\textbackslash end\{longtable\}\\
 definieren\footnote{\url{ftp://ftp.dante.de/pub/tex/macros/latex/required/tools/longtable.pdf}}\\
\begin{longtable}[H]{|p{0.2\textwidth}|p{0.2\textwidth}|p{0.2\textwidth}|}
\hline
A&B&C\\
\hline
\caption{Tabelle 1}
\label{tab:tab1}
\end{longtable}
\subsection{Code-Ausschnitte}
Pseudo-Code Ausschnitte lassen sich mit \textbackslash pseudo\{Name des Algorithmus\}\{Label\}\{Datei im Code-Verzeichnis\} einbinden.
\pseudo{Mittelwert}{lst:mean}{code}
\section{Zitate}
Mit \textbackslash nocite*\{\} lassen sich alle Einträge in der Bibliography ausgeben. Mit \textbackslash cite[S. xx]\{Key\} lassen sich Zitate einfügen. Z.B. \cite[S. 234]{Kurose12} \nocite*{}
\subsection{Abkürzungen}

Abkürzungen können mit \textbackslash nomenclature\{Abk\}\{\textbackslash m\{Abk\}ürzung\} \nomenclature{Abk}{\m{Abk}ürzung} angegeben werden. Diese werden alphabetisch sortiert in ein Abkürzungsverzeichnis aufgenommen.


%
%===============================================================================
% LATEX-Dokument: Literaturverzeichnis
%===============================================================================
%
\newpage
\phantomsection
% Einstellungen f\"{u}r Literaturverzeichnis
\addcontentsline{toc}{section}{\bibname}

\bibliographystyle{IEEEtran}
% argument is your BibTeX string definitions and bibliography database(s)
\bibliography{literature/bib}
% Nutzung von Bibtex:
% hier den bib-file einbinden
%
% GATHER{bibfile.bib}
% \footnotesize
% \bibliography{bibfile}
% ansonsten: bbl als tex Datei einbinden
 %\input{KTR-Seminar-Literatur.tex}
%===============================================================================
% LATEX-Dokument: Literaturverzeichnis
%===============================================================================
%
\end{document}
